%% ============================================================================
%% ch2_problem.tex -- Chapter 2: הגדרת הבעיה, שאלת המחקר והיעדים
%% budget: body pages 4-7 (4 pages)
%% Numbers in 2.1 come from scripts/generate_beauty_benchmarks.py (seed 2026).
%% ============================================================================
\chapter{הגדרת הבעיה, שאלת המחקר והיעדים}
\label{ch:problem}

\begin{managerbox}
הגדרה לקויה של הבעיה בשלב המכין מובילה כמעט בוודאות לכשל בשלבי התכן והמימוש. לכן הפרק מנסח את המצב המצוי במספרים, בלי להזכיר את הפתרון, ומתרגם את הפער למפרט של \num{12} דרישות. לכל דרישה יש ערך יעד ושיטת אימות. דרישה בלי שיטת בדיקה היא משאלה \cite{segal2026booklet}.
\end{managerbox}

הפרק עובר בשרשרת קבועה (איור~\ref{fig:chain}): מהמצב המצוי, דרך הפער הכמותי, אל שאלת המחקר ויעדי התכן, ומשם אל מפרט הדרישות שמשמש חוזה בין הצוות, המנחה והשותף.

\begin{reportfigure}{משרשרת הבעיה אל המפרט: כל חוליה נגזרת מקודמתה}{fig:chain}
\begin{tikzpicture}[node distance=0.55cm,
  b/.style={draw=#1,fill=#1!10,rounded corners=3pt,minimum height=1.25cm,text width=2.35cm,
            align=center,font=\small,line width=0.9pt}]
  \node[b=brandgrey] (s1) {\heb{המצב המצוי}\\[-1pt]\scriptsize\num{$\Delta E_{00}\approx 4.9$}};
  \node[b=brandwarn,left=of s1] (s2) {\heb{הפער}\\[-1pt]\scriptsize\heb{פי \num{2.4} מהסף}};
  \node[b=brandaccent,left=of s2] (s3) {\heb{שאלה ויעדים}\\[-1pt]\scriptsize\enlong{H1 + SMART}};
  \node[b=brandsecondary,left=of s3] (s4) {\heb{מפרט דרישות}\\[-1pt]\scriptsize\en{REQ-xxx-01..12}};
  \foreach \a/\b in {s1/s2,s2/s3,s3/s4} \draw[-Stealth,thick,brandprimary] (\a) -- (\b);
\end{tikzpicture}
\end{reportfigure}

מהאיור עולה שהמפרט אינו רשימת משאלות אלא תוצאה: כל ערך יעד בטבלה~\ref{tab:req} נגזר ממספר שנמדד או חושב במצב המצוי.

\section{ניסוח כמותי של הבעיה}
\label{sec:problem-quant}

\subsection{עיוות צבע בתאורה לא מבוקרת}
מצלמת טלפון ממירה את האור שמגיע מהעור לערכי \en{RGB} דרך שלושה שלבים שאינם בשליטת המשתמשת: מקור האור בחדר, איזון הלבן האוטומטי, ועיבוד התמונה. כדי לכמת את העיוות, הצוות בנה סימולציה (\en{Monte Carlo}) של \num{4,000} תמונות סלפי סינתטיות: גווני עור על פני כל הטווח, תאורה בטמפרטורות צבע של \num{2700--6500 K} עם גוון ירקרק-סגלגל אקראי, חשיפה אוטומטית ורעש חיישן. השגיאה נמדדה במדד \num{$\Delta E_{00}$} (\en{CIEDE2000}), שבו ערך של כ-\num{2.0} הוא בקירוב סף ההבחנה של עין אנושית בין שני גוונים סמוכים \cite{sharma2005ciede2000}. מימוש המדד בסימולציה נבדק מול נתוני הבדיקה שפורסמו עמו.

\begin{reportfigure}{המצב המצוי בסימולציה ($n=\num{4,000}$, נתונים סינתטיים): שגיאת צבע העור, ממוצע ואחוזון \num{95}}{fig:calib}
\begin{tikzpicture}
\begin{axis}[xbar, width=11.6cm, height=3.7cm, bar width=7pt,
  xmin=0, xmax=15.5, xlabel={$\Delta E_{00}$ (skin, lower is better)},
  ytick={0,1}, y dir=reverse, restrict y to domain=-0.5:1.5,
  yticklabels={\heb{מצלמה, ללא תיקון},\heb{איזון לבן גלובלי}},
  yticklabel style={font=\small}, xticklabel style={font=\small},
  legend style={at={(0.98,0.04)},anchor=south east,font=\scriptsize},
  enlarge y limits=0.16, xmajorgrids, grid style={black!10},
  nodes near coords, nodes near coords style={font=\tiny}, every node near coord/.append style={xshift=1pt}]
  \addplot[fill=brandsecondary!75,draw=brandsecondary] table[x=mean,y=idx]{figures/data/calibration_bars.csv};
  \addplot[fill=brandaccent!85,draw=brandaccent!80!black] table[x=p95,y=idx]{figures/data/calibration_bars.csv};
  \draw[brandwarn,dashed,thick] (axis cs:2.0,-0.6) -- (axis cs:2.0,1.6);
  \node[font=\tiny,text=brandwarn,anchor=south west] at (axis cs:2.05,-0.55) {JND $\approx 2.0$};
  \legend{Mean, P95}
\end{axis}
\end{tikzpicture}
\end{reportfigure}

איור~\ref{fig:calib} מסכם את התוצאה. ללא תיקון ייעודי, השגיאה הממוצעת היא \num{$\Delta E_{00} = 4.89$} והאחוזון ה-\num{95} מגיע ל-\num{8.89}. רק \num{8.1\%} מהתמונות נמצאות מתחת לסף ההבחנה. הממצא החשוב יותר הוא שאיזון לבן גלובלי בשיטת "העולם האפור" מחמיר את השגיאה ל-\num{9.77} בממוצע: בתמונת סלפי העור תופס את רוב הפריים, והאלגוריתם "מתקן" את גוון העור עצמו אל הניטרלי. הפער בין המצב המצוי (\num{4.89}) לסף ההבחנה (\num{2.0}) הוא פי \num{2.4} בממוצע, ופי \num{4.4} באחוזון ה-\num{95}.

\subsection{מידות גוף ועלות ההדמיה}
מצב דומה קיים במידות. תמונה דו-ממדית בלי קנה מידה אינה מאפשרת להבחין בין גוף גבוה ורחוק לגוף נמוך וקרוב. לפי אומדן הצוות, שיתוקף בניסוי הבסיס (משימה \en{TSK-1.1}), השגיאה האופיינית בהיקף מותן הנאמד מתמונה כזו היא בסדר גודל של \num{$\pm 4.8$ cm}, יותר ממרווח של מידה אחת בטבלת מידות טיפוסית. בצד הכלכלי, הדמיה גנרטיבית של פריט לבוש עולה \num{\$0.045--0.090} לתמונה בתמחור של שירותי ענן מסחריים (אומדן הצוות). במצב המצוי אין מנגנון שמקשר בין העלות הזו לבין הכנסה, ולכן כל משתמשת פעילה היא הפסד.

\subsection{ההשלכות המעשיות של המצב המצוי}
שלושת הפערים מצטברים לשרשרת אחת. שגיאת צבע של \num{$\Delta E_{00} \approx 5$} היא בערך המרחק בין שני גוונים סמוכים בפלטת מייקאפ טיפוסית, ולכן המלצה על גוון "נכון" מתוך תמונה לא מכוילת נבחרת למעשה בהגרלה בין שניים-שלושה גוונים שכנים. שגיאת היקף של כמה סנטימטרים מתורגמת להמלצת מידה שגויה בכל פעם שהמשתמשת נמצאת קרוב לגבול בין מידות. העלות הגנרטיבית הבלתי-מכוסה מביאה את המפעיל להגביל את מספר ההדמיות, ובכך פוגעת דווקא ברכיב שאמור לבנות אמון. את המחיר משלמות קבוצות שונות מאוד: נערות שקונות לפי טרנד, נשים שמשלבות איפור ולבוש בכל גיל, וגברים בעלי מודעות גבוהה לטיפוח, שכולם נשענים על אותה תמונה לא מכוילת. מבחינה סטטיסטית מדובר בשלוש שגיאות מדידה שמתפשטות זו לתוך זו (\en{Error Propagation}): שגיאת הצבע משפיעה על בחירת הגוון, שגיאת ההיקף על בחירת המידה, ושתיהן על ההסתברות להחזרה.

\subsection{הנחות עבודה ואילוצים}
הניסוח הכמותי נשען על ארבע הנחות, שכל אחת מהן תיבדק בשלב המימוש. (\num{1}) המשתמשת מצלמת במכשיר טלפון מדרגת ביניים, עם זיכרון של \num{4--6 GB} ויחידת עיבוד עצבי (\en{NPU}). (\num{2}) העיניים פתוחות ונראות בתמונה, ואין עדשות מגע צבעוניות. מקרה זה יזוהה כחריג ויטופל במסלול חלופי (\en{RSK-ALG-01} בפרק~\ref{ch:risk}). (\num{3}) קטלוג השותף כולל לכל פריט ערכי צבע ומידות, או שאפשר להפיק אותם מתמונות המוצר. (\num{4}) הסכמה מפורשת של המשתמשת לעיבוד נתונים ביומטריים מתקבלת לפני הצילום הראשון. מכאן נגזרים שלושה אילוצי תכן: התמונה הגולמית אינה עוזבת את המכשיר, זמן הכיול על המכשיר קצר משנייה, וכל רכיב בענן מקבל רק וקטורים מעובדים.

\begin{warningbox}
הבעיה אינה "אין לנו סוכן \en{AI} לביוטי". ניסוח כזה הוא פתרון שמתחפש לבעיה. הבעיה היא שגיאת צבע של פי \num{2.4} מסף ההבחנה, שגיאת היקף שגדולה ממידה אחת, ועלות חישוב שאינה מכוסה.
\end{warningbox}

\section{שאלת המחקר מול יעדי התכן}
\label{sec:rq}

הפרויקט משלב שני מסלולים, ולכל אחד דרך הגדרה אחרת (טבלה~\ref{tab:rq}). המסלול המחקרי עוסק בקשר סיבתי אחד, והמסלול ההנדסי עוסק במערכת שעומדת ביעדים מדידים לפי עקרונות \en{SMART} \cite{doran1981smart}.

\textbf{שאלת המחקר:} האם עיגון צבע ביומטרי בעין, המשלב רישום חד-פעמי בהבזק מסך עם צבע העיניים שהמשתמשת מדווחת כהתפלגות קודמת, מקטין את חציון שגיאת גוון העור בתאורה לא מבוקרת אל מתחת ל-\num{$\Delta E_{00} = 1.8$}, בהשוואה לאיזון הלבן של המצלמה?

\textbf{ההשערות ותכנון הניתוח.} השאלה מפורקת לשתי השערות, שנקבעות מראש ונבחנות עם תיקון הולם (\en{Holm}) לשתי השוואות, ברמת מובהקות כוללת \num{$\alpha = 0.05$}:
$H_{1a}$, שיפור: השגיאה בשיטה המוצעת קטנה מהשגיאה של המצלמה באותה משתתפת ובאותה תאורה. יחידת הניתוח היא המשתתפת: ממוצע ההפרשים על פני ארבע התאורות, במבחן וילקוקסון לזוגות (\en{Wilcoxon signed-rank}) חד-צדדי, ובמקביל מודל מעורב עם אפקט אקראי למשתתפת ואפקט קבוע לתאורה.
$H_{1b}$, סף: חציון השגיאה על פני כל הצילומים קטן מ-\num{1.8}. ההשערה נדחית אם הגבול העליון של רווח סמך חד-צדדי של \num{95\%}, שנבנה בדגימה חוזרת (\en{Bootstrap}) של משתתפות שלמות, קטן מ-\num{1.8}.
האחוזון ה-\num{95} מדווח עם רווח סמך באותה שיטה, ואינו מבחן השערה אלא בדיקת דרישה.
\textbf{עוצמה.} לצורך התכנון נלקח גודל אפקט שמרני של \num{$d = 0.4$}, כרבע מגודל האפקט שנמדד בסימולציה ברמת המשתתפת, כדי לא להישען על אופטימיות של נתונים סינתטיים. לעוצמה של \num{0.9} במבחן חד-צדדי נדרשות \num{56} משתתפות. התוכנית מגייסת \num{120}, \num{40} בכל אחת משלוש קבוצות גוון לפי זווית הטיפוס האישית (\en{ITA}), כדי שגם ניתוח תת-הקבוצות יעמוד בעוצמה הנדרשת ולא תיווצר הטיה לטובת גוונים בהירים \cite{thong2023skin}. בדיקת היתכנות סימולטיבית של השיטה מוצגת בסעיף~\ref{sec:arch}.

\begin{reporttable}{מסלול מחקרי מול מסלול תכן מערכתי בפרויקט}{tab:rq}{P{2.3cm} R R}
\headrow \hc{היבט} & \hc{מסלול מחקרי: כיול ביומטרי} & \hc{מסלול תכן: מערכת \en{AuraTwin AI}} \\ \midrule
המוקד & קשר סיבתי בין עוגן העין לדיוק צבע העור & מערכת שעומדת ב-\num{12} דרישות \\
הניסוח & שאלה ושתי השערות, $H_{1a}$ ו-$H_{1b}$ & יעדי \en{SMART} עם ערך יעד \\
מה משווים & שגיאה מזווגת, אותה משתמשת ואותה תאורה & מדד מול ערך יעד בטבלה~\ref{tab:req} \\
הוכחת הצלחה & דחיית $H_{0a}$ ו-$H_{0b}$ (הולם), \num{$n = 120$} & כל הדרישות עוברות את שיטת האימות \\
תוצר & פרק תוצאות ומאמר כנס & אב-טיפוס בסימולציית שבוע \num{12} \\
\end{reporttable}

הטבלה מראה שהמסלולים משלימים זה את זה: התשובה לשאלת המחקר קובעת אם הדרישה \en{REQ-VIS-01} ברת השגה, ואילו שאר הדרישות אינן תלויות בה ויימדדו בכל מקרה.

\textbf{יעדי התכן לפי \en{SMART}:} (\num{1}) עד סוף שבוע \num{8}, צינור הכיול רץ על המכשיר ועומד ב-\en{REQ-VIS-01} על מדגם הפיתוח של ניסוי הבסיס (\num{30} משתתפות), והאימות המלא על \num{120} משתתפות מסתיים בשבוע \num{11}; (\num{2}) עד סוף שבוע \num{10}, אומדן ההיקפים עומד ב-\num{$\text{MAE} \le 1.2$ cm} מול מדידת סרט; (\num{3}) עד סוף שבוע \num{12}, המערכת המלאה עומדת בזמני התגובה ובדרישת המרווח הכלכלי בבדיקת עומס של \num{200} משתמשות מדומות.

\section{מפרט הדרישות ההנדסי}
\label{sec:req}

מפרט הדרישות בנוי לפי עקרונות תקן \en{ISO/IEC/IEEE 29148} \cite{ieee29148}: לכל דרישה מזהה ייחודי, סיווג (תפקודית, ביצועים, אילוץ או אבטחה), ערך יעד ביחידות מדידות, ושיטת אימות (טבלה~\ref{tab:req}).

\begin{reporttable}[\small]{מפרט הדרישות ההנדסי של \en{AuraTwin AI}}{tab:req}{P{2.3cm} R P{1.3cm} P{2.6cm} R}
\headrow \hc{מזהה} & \hc{הדרישה} & \hc{סיווג} & \hc{ערך יעד} & \hc{שיטת אימות} \\ \midrule
\en{REQ-VIS-01} & כיול גוון העור בתאורה לא מבוקרת & ביצועים & \num{$\tilde{\Delta E}_{00} \le 1.8$}; \num{$P_{95} \le 4.5$} & \num{120} משתתפות, \num{4} תאורות, מול ספקטרופוטומטר \\
\en{REQ-VIS-02} & איתור עוגני העין (לובן וקשתית) & תפקודית & \num{$\ge 97\%$}; \num{$\text{IoU} \ge 0.85$} & \num{500} תמונות מתויגות ידנית \\
\en{REQ-BOD-03} & אומדן היקפי חזה, מותן ואגן & ביצועים & \num{$\text{MAE} \le 1.2$ cm} & \num{60} משתתפות, מדידת סרט כפולה \\
\en{REQ-HND-04} & פילוח ציפורניים וכף יד & ביצועים & \num{$\text{IoU} \ge 0.90$} & \num{300} תמונות מתויגות \\
\en{REQ-LAT-05} & זמן המלצה מקצה לקצה & ביצועים & \num{$P_{90} \le 1.8$ s} & בדיקת עומס ברשת \en{4G} מדומה \\
\en{REQ-GEN-06} & זמן הדמיה גנרטיבית לתמונה & ביצועים & \num{$P_{90} \le 4.5$ s} & מדידה על \en{GPU} בדרגת \en{L4} \\
\en{REQ-GEN-07} & שמירת זהות בהדמיה & ביצועים & \num{$\cos_{\text{ArcFace}} \ge 0.88$} & \num{200} זוגות מקור-הדמיה \\
\en{REQ-GEN-08} & נאמנות הפריט בהדמיה & ביצועים & \num{$\text{LPIPS} \le 0.15$} & מדגם בדיקה מזווג \\
\en{REQ-TRD-09} & עדכניות מנוע הטרנדים & תפקודית & \num{$\le 6$ h} & חותמות זמן ביומן האינדוקס \\
\en{REQ-REC-10} & דיוק המלצה ממומנת & ביצועים & \num{$\text{P@5} \ge 0.85$} & \num{3} סטייליסטיות, \num{$\kappa \ge 0.6$} \\
\en{REQ-SEC-11} & פרטיות ביומטרית & אבטחה & \num{0} תמונות שמורות; מחיקה \num{$\le 60$ s} & ביקורת קוד ויומנים, מבחן חדירה \\
\en{REQ-BIZ-12} & מרווח ליחידת הדמיה & אילוץ & הכנסה \num{$\ge 2.4 \times$} עלות \en{GPU} & לוח עלויות לכל הדמיה \\
\end{reporttable}

\textbf{מקור ערכי היעד.} כל ערך יעד נגזר מאחד משלושה מקורות. \emph{סף תפיסתי:} \en{REQ-VIS-01} נקבעה מתחת לסף ההבחנה של כ-\num{2.0}, בחציון ובאחוזון עליון, כי חוויית המשתמשת נפגעת דווקא בזנב ההתפלגות. \emph{סף מוצר:} \num{$\text{MAE} \le 1.2$ cm} הוא כשליש ממרווח טיפוסי של \num{4 cm} בין מידות. כיוון שגם למדידת סרט יש שגיאה, ערך הייחוס הוא ממוצע של שתי מודדות מיומנות, ומהימנות ביניהן (\en{ICC}) מדווחת לצד התוצאה. \num{97\%} באיתור העין נגזר מכך שכשל איתור מעביר את המשתמשת למסלול כיול חלופי, ויעד המוצר הוא שהמסלול יופעל בפחות מ-\num{3\%} מהמקרים. \num{$\kappa \ge 0.6$} מבטיח שתוויות הסטייליסטיות מהימנות מספיק כדי לשמש קנה מידה. \emph{ייחוס לשיטה קיימת:} יעדי ההדמיה, דמיון \en{ArcFace} לשמירת זהות \cite{deng2019arcface} ו-\en{LPIPS} לדמיון תפיסתי \cite{zhang2018lpips}, הם ערכי התחלה שיעודכנו בשבוע \num{3} מול שחזור מאמר הדגל: היעד הסופי הוא לא יותר מ-\num{0.02} מתחת לביצועי הדגל המשוחזר. יעד הפילוח \num{$\text{IoU} \ge 0.90$} שמרני בכוונה, ו-\num{0.94} מוגדר כיעד שאפתני. \en{REQ-BIZ-12} נגזרת מיעד של מרווח גולמי של כ-\num{58\%} ($1 - 1/2.4$), שמכסה את שאר עלויות הענן, והיא קובעת את תקציב צעדי הדיפוזיה שיתואר בפרק~\ref{ch:design}.

\begin{examplebox}[יישום במערכת: מדרישה לבדיקה]
\en{REQ-VIS-01} תיבדק כך: כל משתתפת מצולמת בארבע תאורות (נורת להט, פלורסנט, אור יום מחלון ולד קר), וגוון העור שלה נמדד במקביל בספקטרופוטומטר. השגיאה מחושבת לכל זוג של תמונה ומדידה, והדרישה עוברת רק אם גם החציון וגם האחוזון ה-\num{95} עומדים בערכי היעד.
\end{examplebox}
